%Je nach dem in welcher Sprache ihr euer Paper schreiben wollt, benutzt bitte entweder den Deutschen-Titel oder den Englischen (einfach aus- bzw. einkommentieren mittels '%')

%Deutsch
\section{Stand des Wissens}\label{sec:related_work}

Aktuelle \acp{sps} werden zunehmend als \ac{iot}-basierte Smart-City-Systeme verstanden. Sie sollen Parksuchverkehr reduzieren, freie Stellplätze schneller auffindbar machen und vorhandene Parkflächen effizienter nutzen. Dafür verbinden sie Sensorik, Kommunikationsstandards, Datenverarbeitung und nutzerseitige Anwendungen zu mehrschichtigen Architekturen. \citet{rajSmartParkingSystems2024} beschreiben moderne \acp{sps} im Kontext von \enquote{Parking 4.0} und unterscheiden physische Sensorebene, Datenübertragung, Verarbeitungs- und Anwendungsebene. Ergänzend geben \citet{alamSurveyIoTDriven2023} einen Überblick über \ac{iot}-getriebene Smart-Parking-Management-Systeme, während \citet{fahimSmartParkingSystems2021} technologische Ansätze, Sensoren, Netzwerktechnologien, User Interfaces, Rechenansätze und Services vergleicht.

Ein wesentlicher Forschungsbereich betrifft die Echtzeiterfassung der Parkplatzbelegung. Dafür werden unterschiedliche Technologien eingesetzt, darunter \ac{ir}- und akustische Sensoren, Kameras sowie Induktionsschleifen \citep{fahimSmartParkingSystems2021}. \citet{perkovicSmartParkingSensors2020} untersuchen die Leistungsfähigkeit üblicher Smart-Parking-Sensoren, Funknetzprotokolle und Optimierungsmöglichkeiten. Neuere Arbeiten rücken Edge-, Fog- und Cloud-Ansätze in den Vordergrund, um Belegungsdaten mit geringer Latenz und skalierbarer Verarbeitung bereitzustellen \citep{alaanzyRealTimeSmart2025,kimRealTimeParkingSpace2025,sarkerSmartParkingSystem2020}. 

Neben der Erfassung freier Stellplätze beschäftigt sich die Forschung mit der Nutzung dieser Daten für Steuerungs-, Routing- und Preisentscheidungen. \citet{hassineDynamicPricingRoute2024} schlagen ein Multi-Agent-System mit dynamischer Preisgestaltung und Routenführung vor, um Parkraumnachfrage und Angebot besser abzustimmen. \citet{icarte-ahumadaMultiAgentSystemParking2025} berücksichtigen bei der Parkplatzzuweisung Kriterien wie Nähe zur Destination, Kosten, Sicherheit, Verkehrsfluss, Fahrzeugeigenschaften und Parkplatzmerkmale. \citet{sarkerSmartParkingSystem2020} verbinden Smart Parking zudem mit dynamischer Preisgestaltung, Edge-Cloud-Computing und LoRa-Kommunikation. \acp{sps} können somit auch betriebliche und verkehrliche Entscheidungen unterstützen.

Für Endnutzer:innen entsteht der Nutzen solcher Systeme nur dann, wenn Informationen während der Fahrt sicher und mit geringer zusätzlicher Interaktion verarbeitet werden können. \citet{oviedo-trespalaciosUnderstandingImpactsMobile2016} zeigen, dass mobile Telefonnutzung während der Fahrt Reaktionszeit, visuelle Aufmerksamkeit, Spurhaltung und Fahrzeugkontrolle negativ beeinflussen kann. \citet{fancelloAnalysisVisualDistraction2025} untersuchen ergänzend visuelle Ablenkung durch Smartphone-Nutzung beim Fahren. Smart-Parking-Informationen sollten daher möglichst ablenkungsarm eingebunden werden.

Für Parkhausbetreiber bleibt entscheidend, ob der erwartete Nutzen die Investitions- und Betriebskosten rechtfertigt. Wirtschaftliche Betrachtungen gewinnen daher an Bedeutung, da Sensorik, Installation, Wartung und Cloud-Dienste laufende Kosten verursachen. \citet{wangEconomicsSmartParking2025} untersuchen wirtschaftliche Effekte von Smart Parking und zeigen, dass datenbasierte Analysen zur Optimierung urbaner Parkraumerlöse beitragen können. Dennoch liegt der Schwerpunkt vieler Arbeiten auf einzelnen technischen Komponenten, Algorithmen oder Services, während ein reproduzierbarer Bewertungsrahmen für konkrete Betreiberentscheidungen weniger stark ausgearbeitet ist.

Zusammenfassend zeigt die Literatur, dass \acp{sps} zahlreiche Einzelaspekte adressieren. Sensorbasierte Systeme erfassen freie Stellplätze, Edge- und Cloud-Architekturen ermöglichen skalierbare Echtzeitverarbeitung, und dynamische Preis- beziehungsweise Routingansätze unterstützen eine effizientere Nutzung vorhandener Parkflächen. Gleichzeitig zeigen Studien zur Fahrerablenkung, dass zusätzliche mobile Interaktion sicherheitskritisch sein kann. Die Forschungslücke liegt daher weniger in der technischen Machbarkeit, sondern in der Verbindung eines realitätsnahen \ac{iot}-Prototyps mit einer nachvollziehbaren Kosten-Nutzen-Bewertung aus Sicht eines Parkhausbetreibers. Drive \& Decide positioniert sich in diesem Zwischenraum, indem ein physischer und virtueller Versuchsaufbau mit einem Kostenmodell für \ac{capex}, \ac{opex}, \ac{tco}, \ac{roi} und Break-Even-Betrachtungen verknüpft wird.

